% Options for packages loaded elsewhere
\PassOptionsToPackage{unicode}{hyperref}
\PassOptionsToPackage{hyphens}{url}
\documentclass[
]{article}
\usepackage{xcolor}
\usepackage[margin=1in]{geometry}
\usepackage{amsmath,amssymb}
\setcounter{secnumdepth}{-\maxdimen} % remove section numbering
\usepackage{iftex}
\ifPDFTeX
  \usepackage[T1]{fontenc}
  \usepackage[utf8]{inputenc}
  \usepackage{textcomp} % provide euro and other symbols
\else % if luatex or xetex
  \usepackage{unicode-math} % this also loads fontspec
  \defaultfontfeatures{Scale=MatchLowercase}
  \defaultfontfeatures[\rmfamily]{Ligatures=TeX,Scale=1}
\fi
\usepackage{lmodern}
\ifPDFTeX\else
  % xetex/luatex font selection
\fi
% Use upquote if available, for straight quotes in verbatim environments
\IfFileExists{upquote.sty}{\usepackage{upquote}}{}
\IfFileExists{microtype.sty}{% use microtype if available
  \usepackage[]{microtype}
  \UseMicrotypeSet[protrusion]{basicmath} % disable protrusion for tt fonts
}{}
\makeatletter
\@ifundefined{KOMAClassName}{% if non-KOMA class
  \IfFileExists{parskip.sty}{%
    \usepackage{parskip}
  }{% else
    \setlength{\parindent}{0pt}
    \setlength{\parskip}{6pt plus 2pt minus 1pt}}
}{% if KOMA class
  \KOMAoptions{parskip=half}}
\makeatother
\usepackage{color}
\usepackage{fancyvrb}
\newcommand{\VerbBar}{|}
\newcommand{\VERB}{\Verb[commandchars=\\\{\}]}
\DefineVerbatimEnvironment{Highlighting}{Verbatim}{commandchars=\\\{\}}
% Add ',fontsize=\small' for more characters per line
\usepackage{framed}
\definecolor{shadecolor}{RGB}{248,248,248}
\newenvironment{Shaded}{\begin{snugshade}}{\end{snugshade}}
\newcommand{\AlertTok}[1]{\textcolor[rgb]{0.94,0.16,0.16}{#1}}
\newcommand{\AnnotationTok}[1]{\textcolor[rgb]{0.56,0.35,0.01}{\textbf{\textit{#1}}}}
\newcommand{\AttributeTok}[1]{\textcolor[rgb]{0.13,0.29,0.53}{#1}}
\newcommand{\BaseNTok}[1]{\textcolor[rgb]{0.00,0.00,0.81}{#1}}
\newcommand{\BuiltInTok}[1]{#1}
\newcommand{\CharTok}[1]{\textcolor[rgb]{0.31,0.60,0.02}{#1}}
\newcommand{\CommentTok}[1]{\textcolor[rgb]{0.56,0.35,0.01}{\textit{#1}}}
\newcommand{\CommentVarTok}[1]{\textcolor[rgb]{0.56,0.35,0.01}{\textbf{\textit{#1}}}}
\newcommand{\ConstantTok}[1]{\textcolor[rgb]{0.56,0.35,0.01}{#1}}
\newcommand{\ControlFlowTok}[1]{\textcolor[rgb]{0.13,0.29,0.53}{\textbf{#1}}}
\newcommand{\DataTypeTok}[1]{\textcolor[rgb]{0.13,0.29,0.53}{#1}}
\newcommand{\DecValTok}[1]{\textcolor[rgb]{0.00,0.00,0.81}{#1}}
\newcommand{\DocumentationTok}[1]{\textcolor[rgb]{0.56,0.35,0.01}{\textbf{\textit{#1}}}}
\newcommand{\ErrorTok}[1]{\textcolor[rgb]{0.64,0.00,0.00}{\textbf{#1}}}
\newcommand{\ExtensionTok}[1]{#1}
\newcommand{\FloatTok}[1]{\textcolor[rgb]{0.00,0.00,0.81}{#1}}
\newcommand{\FunctionTok}[1]{\textcolor[rgb]{0.13,0.29,0.53}{\textbf{#1}}}
\newcommand{\ImportTok}[1]{#1}
\newcommand{\InformationTok}[1]{\textcolor[rgb]{0.56,0.35,0.01}{\textbf{\textit{#1}}}}
\newcommand{\KeywordTok}[1]{\textcolor[rgb]{0.13,0.29,0.53}{\textbf{#1}}}
\newcommand{\NormalTok}[1]{#1}
\newcommand{\OperatorTok}[1]{\textcolor[rgb]{0.81,0.36,0.00}{\textbf{#1}}}
\newcommand{\OtherTok}[1]{\textcolor[rgb]{0.56,0.35,0.01}{#1}}
\newcommand{\PreprocessorTok}[1]{\textcolor[rgb]{0.56,0.35,0.01}{\textit{#1}}}
\newcommand{\RegionMarkerTok}[1]{#1}
\newcommand{\SpecialCharTok}[1]{\textcolor[rgb]{0.81,0.36,0.00}{\textbf{#1}}}
\newcommand{\SpecialStringTok}[1]{\textcolor[rgb]{0.31,0.60,0.02}{#1}}
\newcommand{\StringTok}[1]{\textcolor[rgb]{0.31,0.60,0.02}{#1}}
\newcommand{\VariableTok}[1]{\textcolor[rgb]{0.00,0.00,0.00}{#1}}
\newcommand{\VerbatimStringTok}[1]{\textcolor[rgb]{0.31,0.60,0.02}{#1}}
\newcommand{\WarningTok}[1]{\textcolor[rgb]{0.56,0.35,0.01}{\textbf{\textit{#1}}}}
\usepackage{graphicx}
\makeatletter
\newsavebox\pandoc@box
\newcommand*\pandocbounded[1]{% scales image to fit in text height/width
  \sbox\pandoc@box{#1}%
  \Gscale@div\@tempa{\textheight}{\dimexpr\ht\pandoc@box+\dp\pandoc@box\relax}%
  \Gscale@div\@tempb{\linewidth}{\wd\pandoc@box}%
  \ifdim\@tempb\p@<\@tempa\p@\let\@tempa\@tempb\fi% select the smaller of both
  \ifdim\@tempa\p@<\p@\scalebox{\@tempa}{\usebox\pandoc@box}%
  \else\usebox{\pandoc@box}%
  \fi%
}
% Set default figure placement to htbp
\def\fps@figure{htbp}
\makeatother
\setlength{\emergencystretch}{3em} % prevent overfull lines
\providecommand{\tightlist}{%
  \setlength{\itemsep}{0pt}\setlength{\parskip}{0pt}}
\usepackage{bookmark}
\IfFileExists{xurl.sty}{\usepackage{xurl}}{} % add URL line breaks if available
\urlstyle{same}
\hypersetup{
  pdftitle={Time-dependent Brier{,} IPA{,} external validation},
  hidelinks,
  pdfcreator={LaTeX via pandoc}}

\title{Time-dependent Brier, IPA, external validation}
\author{}
\date{\vspace{-2.5em}}

\begin{document}
\maketitle

\textbf{Testing labs} use the main template: Hypothesis → Visualise →
Assumptions → Conduct → Conclude.

\subsection{Learning objectives}\label{learning-objectives}

\begin{itemize}
\tightlist
\item
  Compute a time-dependent Brier score and the Index of Prediction
  Accuracy (IPA).
\item
  Simulate an external-validation cohort with distribution shift and
  evaluate a fitted model on it.
\item
  Describe why a model can have a high internal AUC and a bad external
  IPA.
\end{itemize}

\subsection{Prerequisites}\label{prerequisites}

Survival ML; calibration.

\subsection{Background}\label{background}

Validation is the point at which a prediction model either earns its
keep or does not. Internal validation --- CV or the bootstrap on the
training data --- controls optimism from the fitting step. External
validation applies the fixed model to an independent cohort and measures
discrimination, calibration, and clinical utility. Time- dependent Brier
scores and IPA (1 − Brier/Brier\_null) summarise the cost of a
probabilistic prediction at a given horizon, combining discrimination
and calibration into a single number per time.

Real biomedical models routinely lose performance on external cohorts.
Distribution shift (different age, comorbidities, measurement protocols)
is the usual reason; miscalibration is the usual failure mode. This lab
simulates both and shows the diagnostic signature each leaves behind.

TRIPOD-AI asks you to report all three --- discrimination, calibration,
clinical utility --- on both internal and external data, and to disclose
the shift between them. A reader who only sees AUC on internal data has
no basis for deploying the model.

\subsection{Setup}\label{setup}

\begin{Shaded}
\begin{Highlighting}[]
\FunctionTok{library}\NormalTok{(tidyverse)}
\FunctionTok{library}\NormalTok{(survival)}
\FunctionTok{library}\NormalTok{(pROC)}
\FunctionTok{set.seed}\NormalTok{(}\DecValTok{42}\NormalTok{)}
\FunctionTok{theme\_set}\NormalTok{(}\FunctionTok{theme\_minimal}\NormalTok{(}\AttributeTok{base\_size =} \DecValTok{12}\NormalTok{))}
\end{Highlighting}
\end{Shaded}

\subsection{1. Hypothesis}\label{hypothesis}

A Cox model fit on a development cohort retains its calibration and
discrimination on an external cohort with shifted covariate
distribution.

\subsection{2. Visualise}\label{visualise}

\begin{Shaded}
\begin{Highlighting}[]
\NormalTok{  x }\OtherTok{\textless{}{-}} \FunctionTok{rnorm}\NormalTok{(n, mu\_x)}
\NormalTok{  lp }\OtherTok{\textless{}{-}}\NormalTok{ beta }\SpecialCharTok{*}\NormalTok{ x}
\NormalTok{  t }\OtherTok{\textless{}{-}} \FunctionTok{rexp}\NormalTok{(n, }\AttributeTok{rate =} \FunctionTok{exp}\NormalTok{(lp }\SpecialCharTok{{-}} \DecValTok{1}\NormalTok{))}
\NormalTok{  c }\OtherTok{\textless{}{-}} \FunctionTok{rexp}\NormalTok{(n, }\AttributeTok{rate =} \FloatTok{0.1}\NormalTok{)}
\NormalTok{  time }\OtherTok{\textless{}{-}} \FunctionTok{pmin}\NormalTok{(t, c)}
\NormalTok{  event }\OtherTok{\textless{}{-}} \FunctionTok{as.integer}\NormalTok{(t }\SpecialCharTok{\textless{}=}\NormalTok{ c)}
  \FunctionTok{tibble}\NormalTok{(}\AttributeTok{x =}\NormalTok{ x, }\AttributeTok{time =}\NormalTok{ time, }\AttributeTok{event =}\NormalTok{ event)}
\ErrorTok{\}}
\NormalTok{dev  }\OtherTok{\textless{}{-}} \FunctionTok{make\_cohort}\NormalTok{(}\DecValTok{500}\NormalTok{, }\AttributeTok{mu\_x =} \DecValTok{0}\NormalTok{)}
\NormalTok{ext  }\OtherTok{\textless{}{-}} \FunctionTok{make\_cohort}\NormalTok{(}\DecValTok{500}\NormalTok{, }\AttributeTok{mu\_x =} \FloatTok{1.0}\NormalTok{)   }\CommentTok{\# distribution shift}
\FunctionTok{bind\_rows}\NormalTok{(dev }\SpecialCharTok{|\textgreater{}} \FunctionTok{mutate}\NormalTok{(}\AttributeTok{cohort =} \StringTok{"dev"}\NormalTok{),}
\NormalTok{          ext }\SpecialCharTok{|\textgreater{}} \FunctionTok{mutate}\NormalTok{(}\AttributeTok{cohort =} \StringTok{"ext"}\NormalTok{)) }\SpecialCharTok{|\textgreater{}}
  \FunctionTok{ggplot}\NormalTok{(}\FunctionTok{aes}\NormalTok{(x, }\AttributeTok{fill =}\NormalTok{ cohort)) }\SpecialCharTok{+}
  \FunctionTok{geom\_histogram}\NormalTok{(}\AttributeTok{bins =} \DecValTok{30}\NormalTok{, }\AttributeTok{alpha =} \FloatTok{0.7}\NormalTok{, }\AttributeTok{position =} \StringTok{"identity"}\NormalTok{)}
\end{Highlighting}
\end{Shaded}

\subsection{3. Assumptions}\label{assumptions}

Independent right censoring; Cox proportional hazards in the development
cohort; no time-varying covariates.

\subsection{4. Conduct}\label{conduct}

Fit on dev; predict survival at a fixed horizon on both cohorts.

\begin{Shaded}
\begin{Highlighting}[]
\NormalTok{horizon }\OtherTok{\textless{}{-}} \DecValTok{3}

\NormalTok{predict\_surv }\OtherTok{\textless{}{-}} \ControlFlowTok{function}\NormalTok{(fit, newdata, t) \{}
\NormalTok{  bh }\OtherTok{\textless{}{-}} \FunctionTok{basehaz}\NormalTok{(fit, }\AttributeTok{centered =} \ConstantTok{FALSE}\NormalTok{)}
\NormalTok{  H0 }\OtherTok{\textless{}{-}} \FunctionTok{approx}\NormalTok{(bh}\SpecialCharTok{$}\NormalTok{time, bh}\SpecialCharTok{$}\NormalTok{hazard, }\AttributeTok{xout =}\NormalTok{ t, }\AttributeTok{rule =} \DecValTok{2}\NormalTok{)}\SpecialCharTok{$}\NormalTok{y}
\NormalTok{  lp }\OtherTok{\textless{}{-}} \FunctionTok{predict}\NormalTok{(fit, }\AttributeTok{newdata =}\NormalTok{ newdata, }\AttributeTok{type =} \StringTok{"lp"}\NormalTok{)}
  \FunctionTok{exp}\NormalTok{(}\SpecialCharTok{{-}}\NormalTok{H0 }\SpecialCharTok{*} \FunctionTok{exp}\NormalTok{(lp))}
\NormalTok{\}}

\NormalTok{s\_dev }\OtherTok{\textless{}{-}} \FunctionTok{predict\_surv}\NormalTok{(fit, dev, horizon)}
\NormalTok{s\_ext }\OtherTok{\textless{}{-}} \FunctionTok{predict\_surv}\NormalTok{(fit, ext, horizon)}

\CommentTok{\# Observed indicator (for simple Brier w/out censoring at horizon)}
\NormalTok{obs\_dev }\OtherTok{\textless{}{-}}\NormalTok{ dev}\SpecialCharTok{$}\NormalTok{event }\SpecialCharTok{==} \DecValTok{1} \SpecialCharTok{\&}\NormalTok{ dev}\SpecialCharTok{$}\NormalTok{time }\SpecialCharTok{\textless{}=}\NormalTok{ horizon}
\NormalTok{obs\_ext }\OtherTok{\textless{}{-}}\NormalTok{ ext}\SpecialCharTok{$}\NormalTok{event }\SpecialCharTok{==} \DecValTok{1} \SpecialCharTok{\&}\NormalTok{ ext}\SpecialCharTok{$}\NormalTok{time }\SpecialCharTok{\textless{}=}\NormalTok{ horizon}

\NormalTok{brier }\OtherTok{\textless{}{-}} \ControlFlowTok{function}\NormalTok{(p\_event, obs) }\FunctionTok{mean}\NormalTok{((p\_event }\SpecialCharTok{{-}}\NormalTok{ obs)}\SpecialCharTok{\^{}}\DecValTok{2}\NormalTok{)}
\NormalTok{brier\_null }\OtherTok{\textless{}{-}} \ControlFlowTok{function}\NormalTok{(obs) }\FunctionTok{mean}\NormalTok{((}\FunctionTok{mean}\NormalTok{(obs) }\SpecialCharTok{{-}}\NormalTok{ obs)}\SpecialCharTok{\^{}}\DecValTok{2}\NormalTok{)}
\NormalTok{b\_dev }\OtherTok{\textless{}{-}} \FunctionTok{brier}\NormalTok{(}\DecValTok{1} \SpecialCharTok{{-}}\NormalTok{ s\_dev, obs\_dev); b0\_dev }\OtherTok{\textless{}{-}} \FunctionTok{brier\_null}\NormalTok{(obs\_dev)}
\NormalTok{b\_ext }\OtherTok{\textless{}{-}} \FunctionTok{brier}\NormalTok{(}\DecValTok{1} \SpecialCharTok{{-}}\NormalTok{ s\_ext, obs\_ext); b0\_ext }\OtherTok{\textless{}{-}} \FunctionTok{brier\_null}\NormalTok{(obs\_ext)}
\NormalTok{ipa\_dev }\OtherTok{\textless{}{-}} \DecValTok{1} \SpecialCharTok{{-}}\NormalTok{ b\_dev }\SpecialCharTok{/}\NormalTok{ b0\_dev}
\NormalTok{ipa\_ext }\OtherTok{\textless{}{-}} \DecValTok{1} \SpecialCharTok{{-}}\NormalTok{ b\_ext }\SpecialCharTok{/}\NormalTok{ b0\_ext}

\NormalTok{auc\_dev }\OtherTok{\textless{}{-}} \FunctionTok{as.numeric}\NormalTok{(}\FunctionTok{auc}\NormalTok{(}\FunctionTok{roc}\NormalTok{(obs\_dev, }\DecValTok{1} \SpecialCharTok{{-}}\NormalTok{ s\_dev, }\AttributeTok{quiet =} \ConstantTok{TRUE}\NormalTok{)))}
\NormalTok{auc\_ext }\OtherTok{\textless{}{-}} \FunctionTok{as.numeric}\NormalTok{(}\FunctionTok{auc}\NormalTok{(}\FunctionTok{roc}\NormalTok{(obs\_ext, }\DecValTok{1} \SpecialCharTok{{-}}\NormalTok{ s\_ext, }\AttributeTok{quiet =} \ConstantTok{TRUE}\NormalTok{)))}

\FunctionTok{tibble}\NormalTok{(}\AttributeTok{cohort =} \FunctionTok{c}\NormalTok{(}\StringTok{"dev"}\NormalTok{, }\StringTok{"ext"}\NormalTok{),}
       \AttributeTok{auc =} \FunctionTok{c}\NormalTok{(auc\_dev, auc\_ext),}
       \AttributeTok{brier =} \FunctionTok{c}\NormalTok{(b\_dev, b\_ext),}
       \AttributeTok{ipa  =} \FunctionTok{c}\NormalTok{(ipa\_dev, ipa\_ext))}
\end{Highlighting}
\end{Shaded}

A calibration plot at the horizon.

\begin{Shaded}
\begin{Highlighting}[]
  \FunctionTok{tibble}\NormalTok{(}\AttributeTok{p =}\NormalTok{ p, }\AttributeTok{y =}\NormalTok{ y, }\AttributeTok{cohort =}\NormalTok{ lab) }\SpecialCharTok{|\textgreater{}}
    \FunctionTok{mutate}\NormalTok{(}\AttributeTok{bin =} \FunctionTok{cut}\NormalTok{(p, }\FunctionTok{quantile}\NormalTok{(p, }\FunctionTok{seq}\NormalTok{(}\DecValTok{0}\NormalTok{, }\DecValTok{1}\NormalTok{, }\AttributeTok{by =} \FloatTok{0.1}\NormalTok{)),}
                     \AttributeTok{include.lowest =} \ConstantTok{TRUE}\NormalTok{)) }\SpecialCharTok{|\textgreater{}}
    \FunctionTok{group\_by}\NormalTok{(cohort, bin) }\SpecialCharTok{|\textgreater{}}
    \FunctionTok{summarise}\NormalTok{(}\AttributeTok{mean\_pred =} \FunctionTok{mean}\NormalTok{(p), }\AttributeTok{obs =} \FunctionTok{mean}\NormalTok{(y), }\AttributeTok{.groups =} \StringTok{"drop"}\NormalTok{)}
\ErrorTok{\}}
\FunctionTok{bind\_rows}\NormalTok{(}\FunctionTok{mk\_cal}\NormalTok{(}\DecValTok{1} \SpecialCharTok{{-}}\NormalTok{ s\_dev, obs\_dev, }\StringTok{"dev"}\NormalTok{),}
          \FunctionTok{mk\_cal}\NormalTok{(}\DecValTok{1} \SpecialCharTok{{-}}\NormalTok{ s\_ext, obs\_ext, }\StringTok{"ext"}\NormalTok{)) }\SpecialCharTok{|\textgreater{}}
  \FunctionTok{ggplot}\NormalTok{(}\FunctionTok{aes}\NormalTok{(mean\_pred, obs, }\AttributeTok{colour =}\NormalTok{ cohort)) }\SpecialCharTok{+}
  \FunctionTok{geom\_point}\NormalTok{() }\SpecialCharTok{+} \FunctionTok{geom\_line}\NormalTok{() }\SpecialCharTok{+}
  \FunctionTok{geom\_abline}\NormalTok{(}\AttributeTok{slope =} \DecValTok{1}\NormalTok{, }\AttributeTok{intercept =} \DecValTok{0}\NormalTok{, }\AttributeTok{colour =} \StringTok{"grey50"}\NormalTok{) }\SpecialCharTok{+}
  \FunctionTok{labs}\NormalTok{(}\AttributeTok{x =} \StringTok{"predicted event prob at horizon"}\NormalTok{,}
       \AttributeTok{y =} \StringTok{"observed event proportion"}\NormalTok{)}
\end{Highlighting}
\end{Shaded}

\subsection{5. Concluding statement}\label{concluding-statement}

\begin{quote}
On a simulated external cohort with shifted covariate distribution, a
Cox model fit on the development cohort retained discrimination (AUC
\texttt{round(auc\_ext,\ 2)}) but lost calibration, with IPA falling
from \texttt{round(ipa\_dev,\ 2)} (development) to
\texttt{round(ipa\_ext,\ 2)} (external). The calibration plot showed
systematic over-prediction in the shifted cohort.
\end{quote}

Distribution shift does not always show up in AUC; it often shows up
first in calibration. Reporting Brier/IPA alongside AUC catches this
failure mode.

\subsection{Common pitfalls}\label{common-pitfalls}

\begin{itemize}
\tightlist
\item
  Reporting AUC as the only discrimination metric on an external cohort.
\item
  Ignoring censoring when computing a Brier score at a horizon;
  \texttt{pec} or \texttt{riskRegression} handle this properly.
\item
  Recalibrating on external data and then reporting the resulting
  performance as external validation.
\end{itemize}

\subsection{Further reading}\label{further-reading}

\begin{itemize}
\tightlist
\item
  Kattan MW, Gerds TA (2018), \emph{The index of prediction accuracy: an
  intuitive measure useful for evaluating risk prediction models}.
\item
  Steyerberg EW, \emph{Clinical Prediction Models}, ch.~15--17.
\end{itemize}

\subsection{Session info}\label{session-info}

\begin{Shaded}
\begin{Highlighting}[]

\end{Highlighting}
\end{Shaded}

\subsection{See also --- chapter index}\label{see-also-chapter-index}

\begin{itemize}
\tightlist
\item
  \href{../index.Rmd}{Methods in this chapter}
\item
  \href{../../../ch00-find-your-method.Rmd}{Find your method (Chapter
  0)}
\end{itemize}

\end{document}
